\section{Resolving the hypergeometric endpoint before assembling the kernel}
\label{sec:endpoint}
We now give an exact integral representation and a two-term asymptotic
for each balanced hypergeometric factor. This is a special-function step
towards \eqref{eq:local-racah-hypothesis}, not a proof of that hypothesis.
The four beta-like prefactors in (B.21)--(B.24) of~\cite{KL} still have
to be combined, particularly on their removable-pole surfaces.

\subsection{An exact Euler representation}
Write a factor, after permuting its parameters, as
\begin{equation}
 H_\Lambda={}_4F_3(A,C,b,d;E,F,g;1),\quad
 A=W+A_0,\quad C=W+C_0,\quad W=\pm2i\Lambda,
 \label{eq:endpoint-H}
\end{equation}
where $E=A+\alpha$, $F=C+\beta$, $g=b+d+\delta$, and
$\alpha+\beta+\delta=1$. All of $A_0,C_0,b,d,\alpha,\beta$
have real part $1/2$ in the real-label problem; $\delta$ is purely
imaginary. The offsets range over a fixed compact set and $\Lambda>0$.
Let $\mathrm B$ denote Euler's beta function, not the function in KL (B.23).
Twice applying the Euler transform~\cite[16.5.2]{DLMF} gives
\begin{align}
 H_\Lambda&=\mathcal N_\Lambda\int_0^1\!\int_0^1
 u^{A-1}(1-u)^{\alpha-1}v^{C-1}(1-v)^{\beta-1}
 {}_2F_1(b,d;g;uv)\,\dd u\dd v,\notag\\
 \mathcal N_\Lambda&=
 \frac{\Gamma(E)\Gamma(F)}
 {\Gamma(A)\Gamma(C)\Gamma(\alpha)\Gamma(\beta)}.
 \label{eq:endpoint-exact-euler}
\end{align}
The identity includes the argument one. One may first insert an argument
$0<z<1$ and pass to $z=1$. The Euler representation of the Gauss function
has a third weight with real exponents $-1/2$, and an additional factor
$(1-zuvt)^{-b}$. Near $u=v=t=1$ its modulus is bounded by a constant times
$[\xi\eta\zeta(\xi+\eta+\zeta)]^{-1/2}$, where
$\xi=1-u$, $\eta=1-v$, $\zeta=1-t$. This is integrable in three dimensions.
All other faces are integrable beta endpoints. Dominated convergence
therefore justifies the endpoint identity without using a divergent
Gauss series at one.

For $\delta\ne0$, the ordinary, unregularized connection formula is
\begin{align}
 {}_2F_1(b,d;g;z)
 &=G_+\,{}_2F_1(b,d;1-\delta;1-z)\notag\\
 &\quad+G_- (1-z)^\delta
 {}_2F_1(g-b,g-d;1+\delta;1-z),\notag\\
 G_+&=\frac{\Gamma(g)\Gamma(\delta)}
 {\Gamma(g-b)\Gamma(g-d)},\qquad
 G_-=\frac{\Gamma(g)\Gamma(-\delta)}{\Gamma(b)\Gamma(d)}.
 \label{eq:endpoint-connection}
\end{align}
This follows from~\cite[15.8.4]{DLMF} after restoring $\Gamma(g)$.
All powers here use the real logarithm for $0<z<1$.

\subsection{A uniform two-branch expansion}
\begin{proposition}
For the compact real-offset families specified above, and any fixed
$0<\epsilon<1/4$,
\begin{equation}
 H_\Lambda=G_++\frac{G_-}{\Gamma(1-\delta)}W^{-\delta}
                 +O(\Lambda^{-1+\epsilon}).
 \label{eq:endpoint-two-branch}
\end{equation}
The error is uniform, including at $\delta=0$, where the sum is interpreted
by its removable limit. With $\log W=\log(2\Lambda)\pm i\pi/2$, that limit is
\begin{equation}
 \frac{\Gamma(b+d)}{\Gamma(b)\Gamma(d)}
       [\log W+\psi(1)-\psi(b)-\psi(d)].
 \label{eq:endpoint-zero}
\end{equation}
No uniform estimate for the complete Racah kernel is asserted here.
\label{prop:endpoint}
\end{proposition}
\begin{proof}
Set $u=e^{-p}$, $v=e^{-q}$, $w=p+q$, $t=p/w$, and
$a=\alpha+\beta=1-\delta$. Equation~\eqref{eq:endpoint-exact-euler} becomes
\begin{equation}
 H_\Lambda=\mathcal N_\Lambda\int_0^\infty e^{-Ww}
        w^{a-1}J(w)\,{}_2F_1(b,d;g;e^{-w})\,\dd w,
 \label{eq:endpoint-fourier}
\end{equation}
where, putting $h(z)=(1-e^{-z})/z$ with $h(0)=1$,
\begin{align}
 J(w)=\int_0^1 &t^{\alpha-1}(1-t)^{\beta-1}
 e^{-w[A_0t+C_0(1-t)]}\notag\\
 &\times h(wt)^{\alpha-1}h(w(1-t))^{\beta-1}\dd t.
 \label{eq:endpoint-J}
\end{align}
This is smooth near zero, uniformly on the compact parameter sets, and
$J(0)=\mathrm B(\alpha,\beta)$. At large $w$ the integrand amplitude
and all its derivatives decay exponentially times a polynomial. This
follows from $\Re A_0,\Re C_0>0$, the behavior $h(z)\sim z^{-1}$,
and the convergent Gauss series at $e^{-w}=0$.

Away from $\delta=0$, the connection formula splits the amplitude at zero
into $G_+w^{a-1}$ and $G_-w^{a+\delta-1}=G_-$, each multiplied by a smooth
factor with value $\mathrm B(\alpha,\beta)$. A smooth cutoff at zero and
the Abel boundary value of the Laplace integral give the leading Fourier
integrals $\Gamma(a)W^{-a}$ and $W^{-1}$. The remainders have powers
$w^a$ and $w$ with bounded differentiated amplitudes. Splitting at
$w=|W|^{-1}$ and integrating by parts on dyadic annuli (more than three
times if needed) bounds their transforms by
$O(|W|^{-\Re a-1})$ and $O(|W|^{-2})$, respectively. The part away from
zero is rapidly decreasing by integration by parts.

Uniform gamma-ratio asymptotics give
\[
 \mathcal N_\Lambda\mathrm B(\alpha,\beta)
 =\frac{W^a}{\Gamma(a)}[1+O(\Lambda^{-1})].
\]
Consequently the two leading terms are those in
\eqref{eq:endpoint-two-branch}. On compact complex parameter sets with
$|\Re\delta|\le\epsilon$ and separated from zero, the error is
$O(\Lambda^{-1+\epsilon})$; the second term's modulus is at most a constant
times $\Lambda^\epsilon$.

To cover coalescence, keep $A_0,C_0,b,d,\alpha$ fixed and vary $\delta$
holomorphically, setting $\beta=1-\alpha-\delta$ and $g=b+d+\delta$.
On $|\delta|=\epsilon$ all Euler endpoint real parts remain positive.
Both $H_\Lambda$ and the combined leading expression are holomorphic
inside this circle. Indeed $\Gamma(-\delta)/\Gamma(1-\delta)=-1/\delta$
makes the cancellation explicit:
\begin{equation}
 G_++\frac{G_-W^{-\delta}}{\Gamma(1-\delta)}
 =\frac{\Gamma(g)}{\Gamma(b)\Gamma(d)}
 \frac{\displaystyle
  \frac{\Gamma(1+\delta)\Gamma(b)\Gamma(d)}
       {\Gamma(b+\delta)\Gamma(d+\delta)}-e^{-\delta\log W}}
 {\delta}.
 \label{eq:endpoint-stable}
\end{equation}
The maximum-modulus principle transfers the boundary error estimate
to the disk. Expanding its numerator at zero gives
\eqref{eq:endpoint-zero}. A finite compact covering completes uniformity.
\end{proof}

For KL (B.33), with $s_i=\Lambda+q_i$ and $s_A=s_B=\nu$, one choice is
\begin{align}
 A_0&=\tfrac12-i(q_3+q_4+\nu),&
 C_0&=\tfrac12-i(q_2+q_3-\nu),\notag\\
 b&=\tfrac12+i(q_4-q_3-\nu),&
 d&=\tfrac12+i(q_2-q_3+\nu),\notag\\
 \alpha&=\tfrac12+i(q_4-q_3+\nu),&
 \beta&=\tfrac12+i(q_2-q_1-\nu),\notag\\
 \delta&=i(q_1-q_2+q_3-q_4),&W&=-2i\Lambda.
 \label{eq:endpoint-B33-map}
\end{align}
The other three factors in (B.21)--(B.24) have the same structure after
permutations, with either sign of $W$. For a fixed nonzero
$\delta=i\chi$, $|W^{-\delta}|=e^{\pm\pi\chi/2}$ is independent of
$\Lambda$. Thus a single hypergeometric factor cannot generally be
replaced pointwise by $G_+$ alone.

This does not refute the assembled distributional asymptotic in KL:
the extra terms may cancel between the four branches. Our finite-point
diagnostics show substantial cancellation at one generic offset tuple,
but do not prove it or cover the opposite-label coincidence surfaces.
The next step is to multiply by the exact beta-like prefactors, prove the
cancellation of the additional leading terms, and bound the combined
remainder across those surfaces before applying the spectral-tail result.
